\documentclass[11pt,a4paper]{article}

% Pacotes essenciais para língua portuguesa e codificação
\usepackage[utf8]{inputenc}
\usepackage[T1]{fontenc}
\usepackage[brazil]{babel}

% Margens e layout
\usepackage[top=2.2cm, bottom=2.2cm, left=2cm, right=2cm]{geometry}
\usepackage{fancyhdr}
\usepackage{microtype}

% Tabelas e formatação matemática
\usepackage{booktabs}
\usepackage{tabularx}
\usepackage{array}
\usepackage{multirow}
\usepackage{amsmath}
\usepackage{amssymb}

% Cores e hiperlinks
\usepackage{xcolor}
\usepackage[colorlinks=true, linkcolor=blue!70!black, urlcolor=blue!70!black, citecolor=blue!70!black]{hyperref}

% Listagem de códigos (Python e SQL)
\usepackage{listings}

% Definição de cores para sintaxe
\definecolor{codegreen}{rgb}{0.13,0.54,0.13}
\definecolor{codegray}{rgb}{0.45,0.45,0.45}
\definecolor{codepurple}{rgb}{0.55,0.14,0.70}
\definecolor{codeblue}{rgb}{0.0,0.3,0.7}
\definecolor{backcolour}{rgb}{0.97,0.97,0.98}

% Configuração global do listings
\lstdefinestyle{codeStyle}{
    backgroundcolor=\color{backcolour},   
    commentstyle=\color{codegray}\itshape,
    keywordstyle=\color{codeblue}\bfseries,
    numberstyle=\tiny\color{codegray},
    stringstyle=\color{codegreen},
    basicstyle=\ttfamily\footnotesize,
    breakatwhitespace=false,         
    breaklines=true,                 
    captionpos=b,                    
    keepspaces=true,                 
    numbers=left,                    
    numbersep=6pt,                  
    showspaces=false,                
    showstringspaces=false,
    showtabs=false,                  
    tabsize=2,
    frame=single,
    rulecolor=\color{black!15}
}
\lstset{style=codeStyle}

% Configuração de cabeçalho e rodapé
\pagestyle{fancy}
\fancyhf{}
\fancyhead[L]{\small\textbf{Laboratório de Banco de Dados | Processo ETL na Prática}}
\fancyhead[R]{\small\textit{Costa dos Coqueiros Turismo}}
\fancyfoot[C]{\thepage}
\renewcommand{\headrulewidth}{0.5pt}
\renewcommand{\footrulewidth}{0.3pt}

\title{\vspace{-1.2cm}\textbf{\Large LABORATÓRIO DE BANCO DE DADOS $|$ PROCESSO ETL NA PRÁTICA}\\[0.3cm]
\large \textbf{Da API pública ao banco de dados}\\[0.2cm]
\normalsize Situação-problema: Costa dos Coqueiros Turismo}
\author{\textbf{Equipe de Engenharia de Dados} \quad$|$\quad \textbf{Valor:} 2,0 pontos}
\date{Novembro de 2026}

\begin{document}

\maketitle
\thispagestyle{fancy}

\noindent\textbf{Repositório do Projeto:} \texttt{lab\_etl/lab\_etl.py} \quad$|$\quad \textbf{Banco de Dados:} \texttt{lab\_etl/lab\_etl.db} (SQLite)\\
\textbf{Período Analisado:} 1º de janeiro a 30 de setembro de 2026 (273 dias corridos).

\vspace{0.3cm}
\hrule
\vspace{0.4cm}

% ==============================================================================
% PARTE 0: PREPARAÇÃO DO AMBIENTE
% ==============================================================================
\section*{Parte 0. Preparação do Ambiente}

\subsection*{Programa Curto de Teste (Selic - Banco Central)}
\begin{lstlisting}[language=Python]
import requests

url = "https://api.bcb.gov.br/dados/serie/bcdata.sgs.432/dados/ultimos/3?formato=json"
response = requests.get(url, timeout=15)

print(f"Código de Status: {response.status_code}")
print(f"JSON Recebido: {response.text}")
\end{lstlisting}

\subsection*{0.1 Saída do teste e significado do código de status 200 \textnormal{[0,03]}}

\noindent\textbf{Saída obtida no console:}
\begin{verbatim}
Código de Status: 200
JSON Recebido: [{"data":"07/12/2026","valor":"13.75"},{"data":"08/12/2026","valor":"13.75"},
                {"data":"09/12/2026","valor":"13.75"}]
\end{verbatim}

\noindent\textbf{O que o código de status 200 significa:}\\
O código \textbf{HTTP 200 (OK)} é a resposta padrão de sucesso do protocolo HTTP. Ele atesta que a solicitação enviada pelo cliente foi recebida, interpretada, aceita e processada com pleno êxito pelo servidor do Banco Central do Brasil, o qual retornou os dados solicitados no corpo da mensagem em formato JSON.

\subsection*{0.2 Versão do Python e do pandas \textnormal{[0,02]}}
\begin{itemize}
    \item \textbf{Comandos executados:} \texttt{python --version} e \texttt{pip show pandas}
    \item \textbf{Versão do Python:} \texttt{Python 3.13.7}
    \item \textbf{Versão do pandas:} \texttt{pandas 3.1.0}
\end{itemize}

% ==============================================================================
% PARTE 1: EXTRAÇÃO
% ==============================================================================
\section*{Parte 1. Extração}

\subsection*{1.1 Inspecione a resposta de cada fonte e complete o quadro \textnormal{[0,04]}}

\begin{table}[h!]
\centering
\small
\begin{tabularx}{\textwidth}{p{2.3cm} X X X}
\toprule
\textbf{Fonte} & \textbf{Onde estão os dados no JSON} & \textbf{Tipo dos valores} & \textbf{Um problema que você já percebe} \\
\midrule
\textbf{IBGE} & Lista JSON na raiz, com um objeto por município contendo estruturas aninhadas. & Misto: inteiros para \texttt{id}, texto para nomes e novos objetos aninhados. & Estrutura hierárquica aninhada (\texttt{microrregiao.mesorregiao.nome}), exigindo achatamento (\textit{json\_normalize}) e renomeação. \\
\addlinespace
\textbf{Banco Central} & Lista JSON na raiz contendo pares de chave-valor (\texttt{data} e \texttt{valor}). & Todos os valores chegam como \textbf{texto (string)}, inclusive cotações e datas. & Dados numéricos e datas chegam como texto; datas no formato brasileiro \texttt{DD/MM/AAAA}; ausência de cotações de dólar em fins de semana e feriados bancários. \\
\addlinespace
\textbf{Open-Meteo} & Objeto JSON na raiz; as medições diárias ficam agrupadas na chave \texttt{daily}. & Datas em texto (\texttt{time}) e variáveis climáticas em números decimais (\texttt{float}). & Dados vêm em listas paralelas separadas; a API não informa o código IBGE do município (exige injeção manual) e a chuva vem em milímetros sem classificação qualitativa. \\
\addlinespace
\textbf{BrasilAPI} & Lista JSON na raiz com um objeto por feriado nacional do ano. & Todos os valores chegam como \textbf{texto (string)} (\texttt{date}, \texttt{name}, \texttt{type}). & Retorna apenas a lista dos dias que são feriados (14 registros), e não o calendário completo contínuo de 273 dias. \\
\bottomrule
\end{tabularx}
\end{table}

\subsection*{1.2 Funções de extração implementadas no script \texttt{lab\_etl.py} \textnormal{[0,12]}}
Foram desenvolvidas 4 funções modulares de extração que salvam a resposta original íntegra na pasta \texttt{raw/} com data e hora no nome e devolvem um DataFrame com todas as colunas convertidas para string (\texttt{astype(str)}):
\begin{enumerate}
    \item \texttt{extrair\_municipios()}: extrai localidades de Sergipe no IBGE, achatando os dados aninhados.
    \item \texttt{extrair\_bcb(serie)}: extrai as séries temporais \texttt{1} (dólar venda) e \texttt{432} (Selic) do Banco Central para o período de 01/01/2026 a 30/09/2026.
    \item \texttt{extrair\_clima()}: extrai temperatura máxima, mínima e precipitação diária de Aracaju no Open-Meteo com timezone \texttt{America/Sao\_Paulo}.
    \item \texttt{extrair\_feriados()}: extrai a lista de feriados nacionais de 2026 da BrasilAPI.
\end{enumerate}

\subsection*{1.3 Função única de chamada HTTP e Espera Exponencial \textnormal{[0,04]}}
A função \texttt{requisicao\_com\_retry(url, params=None, max\_tentativas=3, timeout=30)} possui timeout de 30 segundos, verifica \texttt{status\_code == 200} e realiza até 3 tentativas com pausas progressivas de 1, 2 e 4 segundos.

\vspace{0.2cm}
\noindent\textbf{Por que esperar mais a cada tentativa?}\\
A técnica de \textbf{espera exponencial (\textit{exponential backoff})} é indispensável ao consumir APIs públicas. Falhas temporárias de requisição frequentemente ocorrem por picos de tráfego, instabilidade transitória de rede ou ativação temporária de limite de requisições (\textit{rate limiting}). Tentar novamente de forma imediata ou em laço contínuo sobrecarrega o servidor (gerando uma tempestade de requisições) e pode resultar em bloqueio definitivo do IP. Aumentar progressivamente o intervalo entre as falhas dá fôlego para o servidor remoto recuperar seus recursos e processar a requisição com sucesso.

% ==============================================================================
% PARTE 2: STAGING
% ==============================================================================
\section*{Parte 2. Staging}

\subsection*{2.1 Função de staging \textnormal{[0,04]}}
Implementada na função \texttt{salvar\_staging(conn, df, nome\_tabela, id\_exec)}:
\begin{itemize}
    \item Acrescenta as colunas de auditoria \texttt{id\_exec} (inteiro) e \texttt{dt\_carga} (timestamp da carga).
    \item Substitui pontos e hífens por sublinhado nos nomes das colunas: \texttt{c.replace('.', '\_').replace('-', '\_')}.
    \item Converte todas as colunas para formato texto e salva no banco SQLite \texttt{lab\_etl.db} substituindo a versão anterior (\texttt{if\_exists="replace"}).
\end{itemize}

\subsection*{2.2 Preencha quantas linhas cada fonte entregou \textnormal{[0,02]}}

\begin{table}[h!]
\centering
\small
\begin{tabularx}{0.9\textwidth}{l c X}
\toprule
\textbf{Fonte} & \textbf{Linhas na staging} & \textbf{Comentário} \\
\midrule
\textbf{IBGE, municípios} & \textbf{75} & Total de municípios que formam o estado de Sergipe. \\
\textbf{Banco Central, dólar} & \textbf{187} & Apenas dias úteis de negociação cambial no período (sem fins de semana e feriados bancários). \\
\textbf{Banco Central, Selic} & \textbf{273} & Meta da Selic em vigor cobrindo todos os 273 dias corridos de 01/01 a 30/09/2026. \\
\textbf{Open-Meteo, clima} & \textbf{273} & Série temporal meteorológica contínua para Aracaju em todos os 273 dias. \\
\textbf{BrasilAPI, feriados} & \textbf{14} & Feriados nacionais previstos para 2026 (9 deles situados entre janeiro e setembro). \\
\bottomrule
\end{tabularx}
\end{table}

\subsection*{2.3 Cinco primeiras linhas de \texttt{stg\_dolar} e explicação \textnormal{[0,02]}}

\noindent\textbf{Consulta SQL executada:} \texttt{SELECT * FROM stg\_dolar LIMIT 5;}

\begin{table}[h!]
\centering
\small
\begin{tabular}{cccc}
\toprule
\textbf{data} & \textbf{valor} & \textbf{id\_exec} & \textbf{dt\_carga} \\
\midrule
02/01/2026 & 5.4372 & 2 & 2026-11-03 20:51:37 \\
05/01/2026 & 5.4351 & 2 & 2026-11-03 20:51:37 \\
06/01/2026 & 5.3797 & 2 & 2026-11-03 20:51:37 \\
07/01/2026 & 5.3880 & 2 & 2026-11-03 20:51:37 \\
08/01/2026 & 5.3860 & 2 & 2026-11-03 20:51:37 \\
\bottomrule
\end{tabular}
\end{table}

\noindent\textbf{Por que as colunas da staging são todas texto?}\\
A camada de staging tem como papel primordial armazenar o dado bruto exatamente no estado em que foi transmitido pela fonte externa (\textit{raw ingestion}). Armazenar tudo como texto previne falhas de coerção precoce (\textit{parsing errors}), perda de precisão ou quebras de execução causadas por dados corrompidos ou em formatos inesperados. Isso garante que a ingestão nunca falhe e que qualquer transformação ou limpeza possa ser reexecutada diretamente da staging sem necessidade de novas chamadas de rede.

\subsection*{2.4 Pasta \texttt{raw} \textnormal{[0,02]}}
\noindent\textbf{Arquivos salvos na pasta \texttt{raw/}:}
Foram gerados arquivos JSON com data e hora no nome: \texttt{ibge\_municipios\_*.json}, \texttt{bcb\_dolar\_serie\_1\_*.json}, \texttt{bcb\_serie\_432\_*.json}, \texttt{brasilapi\_feriados\_2026\_*.json}, \texttt{open\_meteo\_clima\_aracaju\_*.json} e \texttt{open\_meteo\_clima\_estancia\_*.json}.

\vspace{0.2cm}
\noindent\textbf{Por que guardar a resposta original da API?}\\
Guardar o payload bruto garante auditoria, histórico e rastreabilidade total do dado de origem. Cria um repositório imutável (\textit{Data Lake local}) que viabiliza a reprodução fiel do pipeline mesmo que as APIs remotas fiquem fora do ar ou alterem sua estrutura de resposta, assegurando resiliência e conformidade com o princípio de idempotência.

% ==============================================================================
% PARTE 3: TRANSFORMAÇÃO
% ==============================================================================
\section*{Parte 3. Transformação}

\subsection*{3.1 Quadro de problemas encontrados e regras de tratamento \textnormal{[0,06]}}

\begin{table}[h!]
\centering
\small
\begin{tabularx}{\textwidth}{c p{5.5cm} X}
\toprule
\textbf{Nº} & \textbf{Problema encontrado} & \textbf{Regra de tratamento} \\
\midrule
1 & Município vem com JSON aninhado, como \texttt{microrregiao.mesorregiao.nome}. & Achatar o JSON aninhado com \texttt{pd.json\_normalize} e renomear o campo para \texttt{mesorregiao}. \\
\addlinespace
2 & O id do município chega como texto na staging. & Converter a coluna para número inteiro: \texttt{cod\_ibge = id.astype(int)}. \\
\addlinespace
3 & A data do Banco Central vem como \texttt{28/09/2026}. & Converter para o padrão ISO \texttt{AAAA-MM-DD}, permitindo ordenação cronológica e junção relacional com as outras tabelas. \\
\addlinespace
4 & O valor do Banco Central vem como texto, por exemplo \texttt{5.1570}. & Converter para o tipo numérico de ponto flutuante (\texttt{float}) com \texttt{pd.to\_numeric()}. \\
\addlinespace
5 & O dólar não tem cotação em fins de semana e feriados bancários. & Gerar o calendário contínuo de 273 dias, cruzar via \texttt{LEFT JOIN} e criar a coluna \texttt{dolar\_ultimo\_valido} utilizando preenchimento adiante (\texttt{ffill}), mantendo \texttt{dolar\_venda} com nulo nos dias sem mercado. \\
\addlinespace
6 & A chuva vem em milímetros, sem classificação. & Aplicar regra condicional de negócio classificando a chuva em: \texttt{sem dado} (se nulo), \texttt{sem chuva} (se 0), \texttt{fraca} ($< 5\text{ mm}$), \texttt{moderada} ($5 \le \text{chuva} < 20\text{ mm}$) e \texttt{forte} ($\ge 20\text{ mm}$). \\
\addlinespace
7 & O clima da API não traz o código do município. & Atribuir manualmente o código oficial do IBGE correspondente a Aracaju (\texttt{2800308}), viabilizando a chave estrangeira para \texttt{dim\_municipio}. \\
\addlinespace
8 & Os feriados vêm como lista só dos dias de feriado. & Efetuar \texttt{LEFT JOIN} a partir do calendário \texttt{dim\_data}, definindo \texttt{feriado = 1} e preenchendo o nome na data coincidente, e \texttt{feriado = 0} nos demais dias. \\
\bottomrule
\end{tabularx}
\end{table}

\subsection*{3.2 Funções de transformação implementadas \textnormal{[0,10]}}
\begin{itemize}
    \item \texttt{transformar\_municipios}: extrai e tipa \texttt{cod\_ibge} (inteiro), \texttt{nome}, \texttt{microrregiao}, \texttt{mesorregiao} e \texttt{uf}.
    \item \texttt{transformar\_bcb}: converte datas para \texttt{AAAA-MM-DD} e valores para \texttt{float}.
    \item \texttt{classificar\_chuva} e \texttt{transformar\_clima}: gera medidas numéricas e categoriza a classe de chuva para Aracaju (\texttt{cod\_ibge = 2800308}).
    \item \texttt{montar\_dim\_data}: monta o calendário de 273 dias com colunas \texttt{data}, \texttt{ano}, \texttt{mes}, \texttt{dia}, \texttt{dia\_semana}, \texttt{fim\_de\_semana} (0 ou 1), \texttt{feriado} (0 ou 1) e \texttt{nome\_feriado}.
    \item \texttt{montar\_fato\_indicador}: une Dólar e Selic cobrindo todos os 273 dias com preenchimento \texttt{ffill}.
\end{itemize}

\subsection*{3.3 Problema de descartar linhas do dólar sem cotação \textnormal{[0,04]}}
Se descartássemos os dias sem cotação, \textbf{perderíamos 86 dias (cerca de 31,5\% de todo o período)}, eliminando todos os sábados, domingos e feriados. Em consultas analíticas que cruzam o clima com os indicadores financeiros (como em junções com \texttt{INNER JOIN}), os eventos meteorológicos dos finais de semana seriam excluídos dos relatórios. Para a agência de turismo \textit{Costa dos Coqueiros}, cujo modelo de negócio depende diretamente da venda de pacotes para fins de semana e feriados prolongados, essa perda distorceria os totais mensais de chuva e inviabilizaria o planejamento de campanhas.

% ==============================================================================
% PARTE 4: CARGA
% ==============================================================================
\section*{Parte 4. Carga}

\subsection*{4.1 Esquema do banco e ativação de chaves estrangeiras \textnormal{[0,04]}}

\begin{lstlisting}[language=SQL]
PRAGMA foreign_keys = ON;

CREATE TABLE dim_municipio (
    cod_ibge INTEGER PRIMARY KEY,
    nome TEXT NOT NULL,
    microrregiao TEXT,
    mesorregiao TEXT,
    uf TEXT
);

CREATE TABLE dim_data (
    data TEXT PRIMARY KEY,
    ano INTEGER NOT NULL,
    mes INTEGER NOT NULL,
    dia INTEGER NOT NULL,
    dia_semana INTEGER NOT NULL,
    fim_de_semana INTEGER NOT NULL,
    feriado INTEGER NOT NULL,
    nome_feriado TEXT
);

CREATE TABLE fato_clima_diario (
    data TEXT NOT NULL,
    cod_ibge INTEGER NOT NULL,
    temp_max REAL,
    temp_min REAL,
    chuva_mm REAL,
    classe_chuva TEXT,
    PRIMARY KEY (data, cod_ibge),
    FOREIGN KEY (data) REFERENCES dim_data(data),
    FOREIGN KEY (cod_ibge) REFERENCES dim_municipio(cod_ibge)
);

CREATE TABLE fato_indicador_diario (
    data TEXT PRIMARY KEY,
    dolar_venda REAL,
    dolar_ultimo_valido REAL,
    selic_meta REAL,
    FOREIGN KEY (data) REFERENCES dim_data(data)
);
\end{lstlisting}

\noindent\textbf{Como você ligou as chaves estrangeiras no SQLite e por que isso importa?}\\
As chaves foram ativadas executando a instrução SQL \texttt{PRAGMA foreign\_keys = ON;} imediatamente após abrir a conexão. Por padrão, o SQLite mantém a checagem de integridade referencial desligada por retrocompatibilidade. Ligá-la é crucial para que o SGBD valide e bloqueie a inserção de registros órfãos nas tabelas de fatos que apontem para datas ou municípios não cadastrados nas dimensões.

\subsection*{4.2 Ordem de carga das tabelas \textnormal{[0,03]}}
A ordem estrita de carga é:
\begin{enumerate}
    \item \textbf{Primeiro as Dimensões:} \texttt{dim\_municipio} e \texttt{dim\_data}.
    \item \textbf{Depois as Fatos:} \texttt{fato\_clima\_diario} e \texttt{fato\_indicador\_diario}.
\end{enumerate}
\textbf{O que aconteceria se invertesse a ordem?}\\
Com as chaves estrangeiras ativadas, tentar carregar as tabelas de fatos antes das dimensões resultaria em erro fatal de violação de integridade (\texttt{FOREIGN KEY constraint failed}), abortando a execução do pipeline.

\subsection*{4.3 Contagem das tabelas finais \textnormal{[0,05]}}

\begin{table}[h!]
\centering
\small
\begin{tabular}{l c l}
\toprule
\textbf{Tabela final} & \textbf{Linhas} & \textbf{Consulta usada} \\
\midrule
\texttt{dim\_municipio} & \textbf{75} & \texttt{SELECT COUNT(*) FROM dim\_municipio;} \\
\texttt{dim\_data} & \textbf{273} & \texttt{SELECT COUNT(*) FROM dim\_data;} \\
\texttt{fato\_clima\_diario} & \textbf{273} & \texttt{SELECT COUNT(*) FROM fato\_clima\_diario;} \\
\texttt{fato\_indicador\_diario} & \textbf{273} & \texttt{SELECT COUNT(*) FROM fato\_indicador\_diario;} \\
\bottomrule
\end{tabular}
\end{table}

\subsection*{4.4 Idempotência da carga \textnormal{[0,03]}}
Ao executar o pipeline pela segunda vez consecutiva, o número de registros permaneceu rigorosamente o mesmo em todas as tabelas (75, 273, 273, 273), sem duplicação de dados nem quebra de chaves primárias. Essa garantia foi obtida através de uma rotina de carga idempotente, que executa uma limpeza ordenada das tabelas antes da inserção em uma transação atômica.

% ==============================================================================
% PARTE 5: CONTROLE E QUALIDADE
% ==============================================================================
\section*{Parte 5. Controle e Qualidade dos Dados}

\subsection*{5.1 Regras de qualidade implementadas \textnormal{[0,03]}}
Implementadas na função \texttt{validar\_qualidade()}:
\begin{enumerate}
    \item \textbf{Completude temporal:} a tabela \texttt{dim\_data} deve conter exatamente 273 registros no intervalo de 01/01/2026 a 30/09/2026.
    \item \textbf{Consistência física de precipitação:} a coluna \texttt{chuva\_mm} não pode conter valores negativos ($\ge 0$).
    \item \textbf{Consistência térmica diária:} a temperatura mínima deve ser obrigatoriamente menor ou igual à máxima (\texttt{temp\_min} $\le$ \texttt{temp\_max}).
    \item \textbf{Unicidade e não-nulidade de chaves primárias:} ausência de duplicatas e nulos em \texttt{data} e na chave composta \texttt{(data, cod\_ibge)}.
    \item \textbf{Sanidade econômica:} a taxa Selic meta deve ser estritamente positiva ($> 0$).
\end{enumerate}

\subsection*{5.2 Tabela \texttt{etl\_controle} e explicação do dólar \textnormal{[0,04]}}

\noindent\textbf{Consulta SQL executada:} \texttt{SELECT id\_exec, fonte, status, linhas\_lidas, linhas\_carregadas, mensagem FROM etl\_controle;}

\begin{table}[h!]
\centering
\footnotesize
\begin{tabularx}{\textwidth}{c p{4.2cm} c c c X}
\toprule
\textbf{id\_exec} & \textbf{fonte} & \textbf{status} & \textbf{linhas\_lidas} & \textbf{linhas\_carregadas} & \textbf{mensagem} \\
\midrule
1 & IBGE Localidades & SUCESSO & 75 & 75 & Extração e staging concluídos. \\
2 & BCB Dólar (Série 1) & SUCESSO & 187 & 273 & Extração e staging concluídos. \\
3 & BCB Selic (Série 432) & SUCESSO & 273 & 273 & Extração e staging concluídos. \\
4 & Open-Meteo Aracaju Clima & SUCESSO & 273 & 273 & Extração e staging concluídos. \\
5 & BrasilAPI Feriados & SUCESSO & 14 & 9 & Extração e staging concluídos. \\
6 & Open-Meteo Estância (Livre) & SUCESSO & 273 & 273 & Fonte livre carregada com sucesso. \\
\bottomrule
\end{tabularx}
\end{table}

\noindent\textbf{Por que as linhas do dólar lidas e carregadas não são 273?}\\
A API do Banco Central forneceu apenas \textbf{187 linhas}, pois a cotação é publicada exclusivamente em dias úteis bancários. No entanto, foram carregadas \textbf{273 linhas} na tabela de fatos porque o modelo analítico requer granularidade diária completa. No processo de transformação, efetuou-se a junção com o calendário contínuo de 273 dias, e a coluna \texttt{dolar\_ultimo\_valido} foi preenchida replicando a última cotação conhecida (\textit{forward fill} / \texttt{ffill}), completando as lacunas de fins de semana e feriados.

\subsection*{5.3 Simulação de erro de propósito \textnormal{[0,03]}}
Ao substituir a URL do IBGE por um endereço inexistente, a chamada HTTP falhou nas 3 tentativas de retry e a tabela \texttt{etl\_controle} registrou com sucesso:
\begin{center}
\small
\texttt{id\_exec: 1 | fonte: IBGE Localidades | status: ERRO | linhas\_lidas: 0 | linhas\_carregadas: 0}\\
\textit{Mensagem: Falha ao consultar a URL após 3 tentativas: https://servicodados.ibge.gov.br/...}
\end{center}
Após o teste, o endereço correto foi restabelecido e o status retornou para \textbf{SUCESSO}.

% ==============================================================================
% PARTE 6: RESPONDENDO AO DIRETOR
% ==============================================================================
\section*{Parte 6. Respondendo ao Diretor (Quadro 1)}

\subsection*{P1: Em quais meses choveu forte em Aracaju e quantos dias foram em cada mês?}
\noindent\textbf{Resposta:} No período de 273 dias em Aracaju, choveu forte ($\ge 20\text{ mm}$) em apenas dois meses: \textbf{março} (1 dia, com 26,8 mm) e \textbf{abril} (1 dia, com 33,6 mm).

\begin{lstlisting}[language=SQL]
SELECT d.ano, d.mes, COUNT(*) AS dias_chuva_forte, ROUND(SUM(c.chuva_mm), 2) AS total_chuva_mm
FROM fato_clima_diario c
JOIN dim_data d ON c.data = d.data
WHERE c.classe_chuva = 'forte'
GROUP BY d.ano, d.mes
ORDER BY d.mes;
\end{lstlisting}

\begin{table}[h!]
\centering
\small
\begin{tabular}{cccc}
\toprule
\textbf{Ano} & \textbf{Mês} & \textbf{Dias com Chuva Forte} & \textbf{Chuva Total nesses Dias (mm)} \\
\midrule
2026 & 3 (Março) & 1 & 26,80 mm \\
2026 & 4 (Abril) & 1 & 33,60 mm \\
\bottomrule
\end{tabular}
\end{table}

\subsection*{P2: Como estava o tempo nos feriados nacionais do período?}
\noindent\textbf{Resposta:} Quase todos os feriados nacionais tiveram tempo bom com chuva fraca e máximas entre 28°C e 30°C. O único feriado com chuva forte foi \textbf{Tiradentes (21/04)}, com 33,6 mm, enquanto \textbf{Corpus Christi (04/06)} registrou chuva moderada (6,9 mm).

\begin{lstlisting}[language=SQL]
SELECT d.data, d.nome_feriado, c.temp_max, c.temp_min, c.chuva_mm, c.classe_chuva
FROM dim_data d
JOIN fato_clima_diario c ON d.data = c.data
WHERE d.feriado = 1
ORDER BY d.data;
\end{lstlisting}

\begin{table}[h!]
\centering
\small
\begin{tabular}{clcccc}
\toprule
\textbf{Data} & \textbf{Feriado} & \textbf{Máx (°C)} & \textbf{Mín (°C)} & \textbf{Chuva (mm)} & \textbf{Classificação} \\
\midrule
2026-01-01 & Confraternização mundial & 30.1 & 25.0 & 1.6 & Fraca \\
2026-02-16 & Carnaval & 29.9 & 24.5 & 1.3 & Fraca \\
2026-02-17 & Carnaval & 29.7 & 24.0 & 2.5 & Fraca \\
2026-04-03 & Sexta-feira Santa & 28.1 & 24.1 & 3.1 & Fraca \\
2026-04-05 & Páscoa & 29.4 & 24.7 & 1.5 & Fraca \\
2026-04-21 & Tiradentes & 27.1 & 23.1 & 33.6 & \textbf{Forte} \\
2026-05-01 & Dia do trabalho & 29.0 & 24.3 & 1.6 & Fraca \\
2026-06-04 & Corpus Christi & 28.7 & 23.7 & 6.9 & \textbf{Moderada} \\
2026-09-07 & Independência do Brasil & 28.9 & 24.5 & 0.6 & Fraca \\
\bottomrule
\end{tabular}
\end{table}

\subsection*{P3: Qual foi o dólar médio de cada mês e quanto choveu no mesmo mês?}
\noindent\textbf{Resposta:} O dólar médio foi mais alto em janeiro (R\$ 5,3457) e mais baixo em maio (R\$ 4,9907). As chuvas concentraram-se entre março e junho, atingindo o pico em \textbf{abril (168,0 mm)}.

\begin{lstlisting}[language=SQL]
SELECT d.ano, d.mes, ROUND(AVG(i.dolar_ultimo_valido), 4) AS dolar_medio, ROUND(SUM(c.chuva_mm), 2) AS chuva_total
FROM dim_data d
JOIN fato_clima_diario c ON d.data = c.data
JOIN fato_indicador_diario i ON d.data = i.data
GROUP BY d.ano, d.mes
ORDER BY d.mes;
\end{lstlisting}

\begin{table}[h!]
\centering
\small
\begin{tabular}{cccc}
\toprule
\textbf{Ano} & \textbf{Mês} & \textbf{Dólar Médio do Mês (R\$)} & \textbf{Chuva Total Acumulada (mm)} \\
\midrule
2026 & 1 (Janeiro) & R\$ 5,3457 & 42,80 mm \\
2026 & 2 (Fevereiro) & R\$ 5,2063 & 87,90 mm \\
2026 & 3 (Março) & R\$ 5,2375 & 124,50 mm \\
2026 & 4 (Abril) & R\$ 5,0381 & 168,00 mm \\
2026 & 5 (Maio) & R\$ 4,9907 & 124,50 mm \\
2026 & 6 (Junho) & R\$ 5,1254 & 151,90 mm \\
2026 & 7 (Julho) & R\$ 5,1145 & 90,40 mm \\
2026 & 8 (Agosto) & R\$ 5,1524 & 74,60 mm \\
2026 & 9 (Setembro) & R\$ 5,1437 & 41,20 mm \\
\bottomrule
\end{tabular}
\end{table}

\subsection*{P4: Em quantos dias a Selic mudou e quais foram a menor e a maior taxa do período?}
\noindent\textbf{Resposta:} A taxa Selic variou de \textbf{13,75\%} (mínima) a \textbf{15,00\%} (máxima), sofrendo alterações em \textbf{5 dias} ao longo de 2026.

\begin{lstlisting}[language=SQL]
WITH mudancas AS (
    SELECT data, selic_meta, LAG(selic_meta) OVER (ORDER BY data) AS selic_anterior
    FROM fato_indicador_diario
)
SELECT ROUND(MIN(selic_meta), 2) AS menor_selic, ROUND(MAX(selic_meta), 2) AS maior_selic,
       COUNT(CASE WHEN selic_meta != selic_anterior THEN 1 END) AS qtd_mudancas
FROM mudancas;
\end{lstlisting}

\begin{table}[h!]
\centering
\small
\begin{tabular}{ccc}
\toprule
\textbf{Menor Selic} & \textbf{Maior Selic} & \textbf{Dias com Mudança de Taxa} \\
\midrule
13,75\% & 15,00\% & 5 \\
\bottomrule
\end{tabular}
\end{table}

\subsection*{P5: Quantos municípios cada mesorregião de Sergipe tem?}
\noindent\textbf{Resposta:} Sergipe possui 75 municípios: 42 no Leste Sergipano (56\%), 18 no Agreste (24\%) e 15 no Sertão (20\%).

\begin{lstlisting}[language=SQL]
SELECT mesorregiao, COUNT(*) AS total_municipios
FROM dim_municipio
GROUP BY mesorregiao
ORDER BY total_municipios DESC;
\end{lstlisting}

\begin{table}[h!]
\centering
\small
\begin{tabular}{lcc}
\toprule
\textbf{Mesorregião} & \textbf{Total de Municípios} & \textbf{Percentual} \\
\midrule
Leste Sergipano & 42 & 56,0\% \\
Agreste Sergipano & 18 & 24,0\% \\
Sertão Sergipano & 15 & 20,0\% \\
\midrule
\textbf{Total Geral} & \textbf{75} & \textbf{100,0\%} \\
\bottomrule
\end{tabular}
\end{table}

% ==============================================================================
% PARTE 7: FONTE LIVRE DA EQUIPE
% ==============================================================================
\section*{Parte 7. Fonte Livre da Equipe}

\subsection*{7.1 Ficha da fonte escolhida \textnormal{[0,10]}}

\begin{table}[h!]
\centering
\small
\begin{tabularx}{\textwidth}{p{4.2cm} X}
\toprule
\textbf{Item} & \textbf{Resposta da equipe} \\
\midrule
\textbf{API escolhida e endereço} & \textbf{Open-Meteo Histórico (Clima de Estância - Litoral Sul de Sergipe)}\newline
\texttt{https://archive-api.open-meteo.com/v1/archive?latitude=-11.26\&longitude=-37.43\&start\_date=2026-01-01\&end\_date=2026-09-30\&daily=temperature\_2m\_max,temperature\_2m\_min,precipitation\_sum\&timezone=America/Sao\_Paulo} \\
\addlinespace
\textbf{Linhas extraídas} & \textbf{273 linhas} (cobrindo diariamente o período de 01/01/2026 a 30/09/2026). \\
\addlinespace
\textbf{Duas regras de transformação aplicadas} & 1. \textbf{Injeção do código IBGE:} Associação da chave primária de Estância (\texttt{cod\_ibge = 2802106}) para integridade referencial com \texttt{dim\_municipio}.\newline
2. \textbf{Classificação qualitativa de precipitação:} Categorização condicional nas faixas (\texttt{sem dado}, \texttt{sem chuva}, \texttt{fraca}, \texttt{moderada}, \texttt{forte}). \\
\addlinespace
\textbf{Tabela final criada, com a chave primária} & \textbf{\texttt{fato\_clima\_estancia}}\newline
Chave primária: \texttt{data TEXT PRIMARY KEY}, com chaves estrangeiras apontando para \texttt{dim\_data(data)} e \texttt{dim\_municipio(cod\_ibge)}. \\
\addlinespace
\textbf{Pergunta de negócio e consulta SQL} & \textbf{Pergunta:} \textit{Quais foram os melhores fins de semana de sol (sem chuva e temperatura máxima $\ge 29^\circ\text{C}$) no litoral sul (Estância/Praia do Saco) para direcionar as campanhas de pacotes turísticos?}\newline\newline
\textbf{Consulta SQL:}\newline
\texttt{SELECT d.data, CASE d.dia\_semana WHEN 5 THEN 'Sábado' WHEN 6 THEN 'Domingo' END AS dia\_fim\_semana, e.temp\_max, e.chuva\_mm, e.classe\_chuva FROM fato\_clima\_estancia e JOIN dim\_data d ON e.data = d.data WHERE d.fim\_de\_semana = 1 AND e.chuva\_mm = 0 AND e.temp\_max >= 29.0 ORDER BY d.data;} \\
\bottomrule
\end{tabularx}
\end{table}

\end{document}

